Indianapolis 500 qualifying creates a distinctive sequential decision problem once a car has established a
valid four-lap result: a team may withdraw that result and enter a higher-priority queue in pursuit of an
earlier reattempt, or retain it and wait for another opportunity while the physical environment evolves.
Historical investigation showed that queue entry, lane state, withdrawal timing, requeue behaviour and team
intent could not be reconstructed consistently enough across seasons to support a defensible waiting-time
model. The problem was therefore reformulated conditionally: if another on-track opportunity occurs $h$
minutes from now, what distribution of four-lap performance change should be expected relative to the
car's current official result? The resulting scientific target is $p(\Delta v \mid H=h)$, deliberately not
$P(H=h)$.

A frozen 2020--2024 performance core of 41 evidence-qualified same-car transitions was reconstructed with
explicit provenance, confidence and quarantine rules. A zero-intercept physical-response model links changes
in track and ambient temperature to changes in four-lap average speed; a separately validated future
track-state model propagates horizon, future ambient-temperature change, current thermal gap and solar
state. Three uncertainty sources -- future track-state uncertainty, bootstrap coefficient uncertainty, and
empirical attempt-level residual variation -- are combined by Monte Carlo simulation into calibrated
probabilistic outlooks at five validated horizons (15, 30, 60, 90, 120 minutes). Externally, on 15 eligible
2025 hybrid-era repeat attempts using realised environmental conditions, the frozen inference core achieved
MAE $=0.492$ mph, median absolute error $=0.425$ mph, RMSE $=0.610$ mph, 80\% predictive-interval
coverage of 80\%, 90\% coverage of 100\%, and directional accuracy of 46.7\%.

That evaluation isolates the physical-inference layer using realised future weather rather than a live forecast.
A subsequent point-in-time evaluation layer was therefore constructed on top of the same frozen scientific
core: a forecast-vintage store enforcing \texttt{forecast.issue\_time} $\le$ \texttt{decision\_time}, an
immutable-snapshot historical shadow replay, and evidence-aware abstention that reports, rather than hides,
every case the available evidence cannot support. Of 41 historical same-car transitions, only 10 yield a
genuine, non-fabricated point-in-time decision case; of those, 9 support a full five-horizon conditional
physical outlook while their realised recurrence falls outside the 120-minute calibrated boundary, and exactly
one -- retained explicitly as an illustrative diagnostic, never as validation evidence -- falls near a calibrated
anchor. We report plainly that the point-in-time architecture and leakage controls are demonstrated in
running software, while statistically meaningful point-in-time forecast validation remains, on the currently
available evidence, not established. The system does not predict queue waiting time, next-opportunity time,
or an optimal retain/withdraw strategy; it is intended as one quantitative layer within a broader human
pit-wall decision.
